\documentclass[12pt,a4paper]{article}
\usepackage[utf8]{inputenc}
\usepackage[german]{babel}
\usepackage[T1]{fontenc}
\usepackage{geometry}
\usepackage{setspace}
\usepackage{titlesec}
\usepackage{url}
\usepackage{amsmath}
\usepackage{graphicx}
\usepackage{multicol}
\usepackage{datetime}

\geometry{margin=2.5cm}
\onehalfspacing

\titleformat{\section}{\large\bfseries}{\thesection}{1em}{}
\titleformat{\subsection}{\normalsize\bfseries}{\thesubsection}{1em}{}
\setlength{\columnsep}{25pt}

\newdateformat{germandate}{\THEDAY.\THEMONTH.\THEYEAR}

\title{\textbf{Selbstüberwachtes Training für die Datierung von Dokumentbildern}}
\author{Dmytro Pahuba\\
\texttt{dmytro.pahuba@tu-dortmund.de}}
\date{\germandate\today}

\begin{document}

\maketitle

\section*{\centering Zusammenfassung}

In dieser Arbeit wird die Anwendung selbstüberwachter Lernverfahren für die automatisierte Datierung historischer Dokumentenbilder untersucht. Während traditionelle überwachte Methoden durch begrenzte annotierte Datensätze eingeschränkt sind, ermöglicht selbstüberwachtes Vortraining die Nutzung großer Mengen nicht-annotierter historischer Manuskripte. Basierend auf dem AttMask-Framework von Kakogeorgiou et al. \cite{kakogeorgiou2022hide}, das Raven et al. \cite{raven2024self} erfolgreich für Writer Retrieval einsetzten ,  testet diese Arbeit experimentell, ob und inwiefern Merkmale, die zur Unterscheidung individueller Schreiber gelernt wurden, auch für die temporale Klassifikation historischer Dokumente geeignet sind.

\begin{multicols}{2}

\section{Motivation}

Die automatische Datierung historischer Dokumente besitzt hohe Relevanz für die Geschichtsforschung, da sie undatierte Manuskripte kategorisieren und in ihren kulturellen Kontext einordnen kann. Traditionell ist dies das Feld von Paläographen mit jahrelanger Ausbildung, deren manuelle Datierungen jedoch zeitaufwändig und kostspielig sind. Zudem variieren Experteneinschätzungen oft um 50-100 Jahre, wobei ±25 Jahre als akzeptabler Fehlerbereich gelten. Ein zentrales Problem überwachter Lernverfahren ist die begrenzte Verfügbarkeit annotierter Datensätze. Der KERTAS-Datensatz umfasst etwa 2000 arabische Manuskripte \cite{adam2018kertas}. Für datenhungrige Deep Learning Modelle sind das problematisch kleine Zahlen. Selbstüberwachtes Lernen bietet eine Lösung, da es große Mengen nicht-annotierter Manuskripte nutzen kann. Im Writer Retrieval erreichten Raven et al. \cite{raven2024self} mit dem AttMask-Framework 83.1 Prozent mAP auf Historical-WI und 95.0 Prozent mAP auf HisIR19. Da Schreibstile sich über Zeiträume entwickeln, könnten Features zur individuellen Schreiberdifferenzierung auch temporale Klassifikation unterstützen. Allerdings zeigten He et al. \cite{he2016image}, dass Writer ID Features nicht automatisch für Dating geeignet sind, da Dating allgemeine Periodenstile erfasst, während Writer ID individuelle Merkmale fokussiert. Die Übertragbarkeit bleibt daher experimentell zu validieren.

\section{Verwandte Arbeiten}

\subsection{Traditionelle Dokumentdatierung}

Frühe Ansätze nutzten handcrafted Features. He et al. \cite{he2016image} erreichten mit Kontur- und Strichfragmenten 7.9 Jahre MAE auf dem Medieval Paleographic Scale Datensatz. \cite{wahlberg2016historical} erzielten mit CNNs 18.9 Jahre MAE auf schwedischen Urkunden. Diese Ergebnisse sind bereits sehr gut, da Paläographen typischerweise um ±25 Jahre variieren.

\subsection{Deep Learning Ansätze}

Hamid et al. \cite{hamid2019deep} verbesserten die Ergebnisse mit Transfer Learning auf vortrainierten CNNs und erreichten 10.23-16.16 Jahre MAE auf dem MPS-Datensatz . Der ICDAR-2021 Competition Benchmark \cite{seuret2021icdar} setzte den aktuellen Standard mit 22 Jahre MAE als beste Leistung. Aktuelle Vision Transformer Ansätze wie DiT \cite{li2022dit} fokussieren auf moderne Dokumente, nicht historische Manuskripte.

\subsection{Writer Identification und SSL}

Raven et al. \cite{raven2024self} demonstrierten erfolgreiche SSL-Anwendung für Writer Retrieval mit 83.1 Prozent mAP auf Historical-WI. Allerdings ist fraglich, ob Writer-ID-Features für Dating übertragbar sind. He et al. \cite{he2016image} zeigten bereits, dass Features für individuelle Schreiberdifferenzierung schlecht für temporale Stilklassifikation funktionieren, da verschiedene Abstraktionsebenen betroffen sind.

\section{Methodik}

\subsection{Datenaufbereitung}

Die Arbeit nutzt den CLAMM ICDAR-2017-Datensatz \cite{cloppet2017icdar2017} mit 6,540+ Bildern aus 1,100 Jahren (500-1600 CE) und 12 paläographischen Schriftklassen. Zusätzlich wird der CATMuS Medieval-Datensatz eingesetzt. Vorverarbeitung umfasst Normalisierung auf 256$\times$256 Pixel, Graustufenkonvertierung und Augmen-tierung durch Rotation, Skalierung und Rauscheinführung. Kritisch ist die tem-porale Aufteilung der Datensätze zur Vermeidung von Informationsleckage. 

\subsection{SSL-Architektur}

Das System nutzt das AttMask-Framework, das DINO Self-Distillation mit attention-guided Masking kombiniert. Der Teacher identifiziert wichtige Bildregionen über Self-Attention Maps, die dann beim Student maskiert werden. Features werden über VLAD-Aggregation von Foreground-Patches zu globalen Dokumentrepräsentationen zusammengefasst. Die SSL-basierten Features sind zunächst task-agnostisch und lernen allgemeine visuelle Muster in Dokumenten. Erst durch nachgelagertes Training auf Writer Identification werden diese generischen Repräsentationen zu schreiberspezifischen Features spezialisiert. Die zentrale Forschungsfrage ist, ob diese spezialisierten Writer-ID-Features auf Document Dating übertragbar sind, oder ob der Transfer bereits auf der Ebene der generischen SSL-Features erfolgen sollte.

\subsection{Transfer Learning Ansatz}

Die Arbeit untersucht zwei Transfer Learning Strategien für die Nutzung selbstüberwachter Repräsentationen:

\textbf{Strategie 1 - Domain-spezifisches SSL:} AttMask wird direkt auf den Dating-Datensätzen (CLAMM, CATMuS) vortrainiert und anschließend für Dating fine-getunet. Dies nutzt die SSL-Architektur ohne vorherige Task-Spezialisierung.

\textbf{Strategie 2 - Task-Transfer:} Das auf ICDAR2017 Writer-ID vortrainierte Modell wird für Dating adaptiert. Dies testet die Übertragbarkeit bereits spezialisierter Writer-ID-Features auf temporale Klassifikation.

Beide Strategien werden durch Feature Evaluation ohne task-spezifisches Training bewertet:

\begin{enumerate}
\item \textbf{Lineare Evaluation:} Ein linearer Klassifikator auf gefrorenen Features evaluiert die Repräsentationsqualität direkt
\item \textbf{k-NN Evaluation:} k-Nearest-Neighbor Klassifikation mit auf Validierungsdaten optimiertem k testet die Clusterbildung der Features
\end{enumerate}

Diese Evaluationen messen die inhärente Güte der vortrainierten Features für Dating ohne weitere Anpassung. Anschließend erfolgt Fine-tuning des gesamten Modells mit kombiniertem Writer-ID und Dating Loss bei Multi-Task Learning.

\subsection{Evaluation}

Hauptmetrik ist Mean Absolute Error (MAE) in Jahren
% \begin{equation}
% MAE = \frac{1}{n} \sum_{i=1}^{n} |y_i - \hat{y}_i|
% \end{equation}

Zusätzliche Metriken umfassen Root Mean Square Error (RMSE) und Accuracy@k (Vorhersage innerhalb $k$ Jahre). Validierung erfolgt durch temporale Kreuzvalidierung zur Vermeidung von Data Leckage.

\section{Experimente}

\subsection{Datensätze}

\begin{itemize}
\item \textbf{CLAMM ICDAR-2017:} 6,540+ Bilder, 1,100 Jahre Zeitspanne (500-1600 CE), 12 paläographische Schriftklassen
\item \textbf{CATMuS Medieval \cite{clerice2024catmus}:} 200+ Manuskripte und Inkunabeln in 10 Sprachen, 160,000+ Textzeilen, 5 Millionen Zeichen, 8.-16. Jahrhundert mit reichen Metadaten (Jahrhundert, Sprache, Genre, Schrift)
\item \textbf{ICDAR-2021 HDC \cite{seuret2021icdar}:} Competition-Datensatz für historische Dokumentklassifikation mit Dating-Task, beste Systeme erreichten 22 Jahre Mean Absolute Error
\end{itemize}

\subsection{SSL-Pretraining}

Vision Transformer wird mittels AttMask-Framework auf unlabeled historischen Dokumenten vortrainiert. Das SSL-Training lernt generische visuelle Repräsentationen ohne task-spezifische Labels. Die Hyperparameter werden experimentell für jeden Datensatz optimiert.

\subsection{Transfer Learning Pipeline}

\begin{enumerate}
\item SSL-Pretraining auf unlabeled Dokumentenbildern
\item Lineare Evaluation der extrahierten Features
\item k-NN Evaluation ohne Hyperparameter-Tuning
\item Fine-tuning für Writer-ID zu Dating Transfer
\item Vergleich mit überwachter ResNet-50 Baseline
\end{enumerate}

\subsection{Ablationsstudien}

\begin{itemize}
\item Einfluss verschiedener Masking-Ratios (40-75\%)
\item VLAD-Clusterzahl (50, 100, 200)
\item Foreground-Threshold für Patch-Selektion
\item Vergleich flat encoding vs. hierarchische Ansätze
\item Augmentierungsstrategien: geometrische vs. morphologische Transformationen
\end{itemize}


\bibliographystyle{plain}
\bibliography{references}

\end{multicols}

\end{document}